\section{Relation to Prior Work and Positioning}\label{sec:prior-work}

This paper sits at the junction of several established literatures. Its main theorem is new in its exact form, but its proof language and case studies draw on standard work in logic and foundations. The point of this section is therefore not to claim a single narrow ancestry. It is to make clear which literatures the manuscript enters and what it takes from each of them.

\paragraph*{Conservative extension and theory comparison.}
One strand of the paper concerns when an apparent strengthening says nothing genuinely new about the original base. That issue belongs to the familiar study of definitions, expansions, interpretations, and conservative extensions in model theory and adjacent work on theory comparison \citep{ChangKeisler1990,Hodges1993}. More recent discussions of definitional, categorical, and Morita equivalence sharpen the idea that two presentations may differ in signature while agreeing on what is added about their original subject matter \citep{BarrettHalvorson2016}. The present manuscript does not collapse those notions into one equivalence relation. It uses them as anchors for the narrower claim that internally fixed auxiliary data yield only conservative change.

\paragraph*{Partial structures and genuine undefinedness.}
The no-totalization result belongs formally to the setting of partial structures and partial algebras, where undefinedness is part of the structure rather than a temporary defect. Burmeister provides the relevant background vocabulary \citep{Burmeister1986}. The present use of partiality is sharply focused: it makes precise the claim that if an old primitive acquires a new value on an old tuple, then the old base itself has changed.

\paragraph*{External methods in model theory and set theory.}
Several of the paper's main examples are standard external methods of logic. Model theory uses compactness, saturation, monster models, and witness structures to move to richer environments \citep{Hodges1993,Marker2002}. Set theory uses forcing, generic extensions, and absoluteness to compare what survives passage to new universes \citep{Jech2003,Kanamori2008}. The paper does not dispute those practices. It reclassifies their fixed-base status: relative to a fixed base, they are ambient enlargements, witness imports, or conservative readbacks, not autonomous sources of new internal operator power.

\paragraph*{Universal constructions and ambient comparison.}
Category-theoretic existence by universal property is another natural neighboring literature. Universal arrows, adjunctions, limits, colimits, and free constructions are posed relative to an ambient category and comparison maps \citep{MacLane1998}. For the present paper, that ambient dependence is not an objection; it is the feature that marks the boundary between internal generation on an object alone and existence mediated by a wider framework.

\paragraph*{Foundational methodology.}
Methodologically, the closest analogue is fixed-base analysis in foundations. Reverse mathematics asks which set-existence principles are needed over a chosen base theory to prove a theorem \citep{Simpson2009}. The present paper asks a different but structurally similar question: once a base structure is fixed, which operator moves are already internal to it, and which rely on additional witnesses or scope extension? In that sense the manuscript is best read as a boundary theorem about internality rather than as an elimination program.

\paragraph*{Terminological caution and novelty.}
The manuscript's use of ``admissibility'' is local and regime-relative; it is not meant to place the paper inside admissible set theory unless that connection is made explicitly for some later application \citep{Barwise1975}. The novelty claim is likewise specific. The paper is not merely putting conservative extension, forcing, saturation, and universal constructions into one list. It argues that these apparently different families are governed by a common dichotomy: once the load-bearing datum is isolated, the move is either conservative because that datum was already fixed by the base, or external because it was not. Coupled with the rigidity theorem for partial structures and the exhaustion theorem for apparent strengthening mechanisms, that is the unified contribution of the manuscript.
